\begin{lemma}[2.5.6]
\label{IV.2.5.6-fr}
Sean $B$ un anillo no necesariamente conmutativo y $V$, $W$ dos
$B$-módulos a izquierda isomorfos; sea $C$ el anillo
$\operatorname{End}_B(V)$, y sea
$M=\operatorname{Hom}_B(V,W)$, dotado de su estructura canónica de
$C$-módulo a derecha. Entonces $M$ es un $C$-módulo isomorfo a $C_d$; además,
para todo $u\in M$, las condiciones siguientes son equivalentes:
\begin{enumerate}
\item[\emph{a)}] $\{u\}$ es una base del $C$-módulo $M$.
\item[\emph{b)}] $u$ es un isomorfismo de $V$ sobre $W$.
\end{enumerate}
\end{lemma}

Si $u$ es un isomorfismo de $V$ sobre $W$, la aplicación
$v\leadsto u\circ v$ de $C$ en $M$ es evidentemente una biyección, por lo que
\emph{b)} implica \emph{a)}. Recíprocamente, supongamos que $\{u\}$ es una
base del $C$-módulo $M$. Por hipótesis, existe un isomorfismo $u'$ de $V$
sobre $W$, y $\{u'\}$ es, por tanto, una base de $M$; existe entonces un
elemento invertible $w$ de $C$ (es decir, un automorfismo de $V$) tal que
$u=u'\circ w$, lo que implica que $u$ es un isomorfismo de $V$ sobre $W$.

\begin{corollary}[2.5.7]
\label{IV.2.5.7-fr}
Siendo las hipótesis sobre $B$, $V$ y $W$ las de \textup{(2.5.6)}, supongamos además
que se cumple una de las condiciones siguientes:
\begin{enumerate}
\item[\textup{(i)}] $V$ y $W$ son $B$-módulos noetherianos;
\item[\textup{(ii)}] $V$ y $W$ son módulos de presentación finita sobre un subanillo
conmutativo de $B$.
\end{enumerate}
Entonces las condiciones \emph{a)} y \emph{b)} de \textup{(2.5.6)} son también
equivalentes a las siguientes:
\begin{enumerate}
\item[\emph{a}$'$)] $u$ es un generador del $C$-módulo $M$.
\item[\emph{b}$'$)] $u$ es un epimorfismo de $V$ sobre $W$.
\end{enumerate}
\end{corollary}

Se sabe, en efecto, que un epimorfismo de un $A$-módulo $E$ sobre sí mismo es
\emph{biyectivo} en los dos casos siguientes: $1^\circ$ $E$ es un $A$-módulo
noetheriano (Bourbaki, \emph{Alg.}, cap.~VIII, \S~2,
núm.~2, lema~3); $2^\circ$ $A$ es conmutativo y $E$ es un
$A$-módulo de presentación finita (8.9.3)\footnote{El lector podrá comprobar
que (2.5.7) y
(2.5.8) no se utilizan antes del \S~9.}; por tanto, \emph{b)} y
\emph{b}$'$) son equivalentes. Por otra parte, si $u$ genera $M$
\oldpage[IV-2]{25}
y si $\{u'\}$ es una base de $M$, existe $v\in C$ tal que
$u'=u\circ v$, lo que demuestra que $u$ es sobreyectivo; por tanto,
\emph{a}$'$) implica \emph{b}$'$), y, como \emph{a)} implica evidentemente
\emph{a}$'$), esto termina de demostrar el corolario.
